
\subsubsection{LLM Calibration}

Neural network calibration has been studied since Guo et al. (2017) established Expected Calibration Error as the standard metric and introduced temperature scaling. For LLMs, Kadavath et al. (2022) demonstrated that models possess intrinsic self-knowledge and can predict answer correctness at above-chance rates.

Subsequent work explored various elicitation approaches. Lin et al. (2022) trained models to express uncertainty in natural language. Xiong et al. (2023) evaluated confidence elicitation methods including direct prompting, multi-sample consistency, and linguistic confidence markers. Tian et al. (2023) studied prompting strategies for calibration.

\subsubsection{Chain-of-Thought Prompting}

Wei et al. (2022) introduced chain-of-thought prompting for reasoning tasks. Kojima et al. (2022) extended this to zero-shot settings with ''Let's think step by step.'' Wang et al. (2023) introduced self-consistency using multiple samples with majority voting.

These works focus on accuracy rather than calibration. A key observation is that CoT produces visible reasoning traces that could reveal model uncertainty through linguistic markers.

\subsubsection{Uncertainty Quantification in NLP}

Linguistic hedging has long been recognized as marking epistemic uncertainty \citep{hyland1998hedging}. Kuhn et al. (2023) proposed semantic uncertainty, clustering semantically equivalent responses to estimate uncertainty.

The present study differs from prior work by systematically isolating the mechanism by which CoT might improve calibration, decomposing the process into verifiable sub-hypotheses rather than measuring only end-to-end effects.
